% LaTeX Template for Including PNG Figures in arXiv Submission
% Based on arXiv submission guidelines: https://info.arxiv.org/help/submit/index.html

\documentclass[12pt]{article}
\usepackage[utf8]{inputenc}
\usepackage{graphicx}      % Required for \includegraphics
\usepackage{float}         % For [H] figure placement
\usepackage{subcaption}    % For subfigures
\usepackage{booktabs}      % For nice tables
\usepackage{amsmath}       % For math

% ============================================
% FIGURE INCLUSION EXAMPLES
% ============================================

% Example 1: Single figure, full width
\begin{figure}[H]
    \centering
    \includegraphics[width=\textwidth]{figures/temperature_effects.png}
    \caption{Mean alignment scores across temperature range (0.0-1.0) for each model. Error bars represent standard error of the mean (SEM).}
    \label{fig:temperature}
\end{figure}

% Example 2: Single figure, 80% width
\begin{figure}[H]
    \centering
    \includegraphics[width=0.8\textwidth]{figures/framing_effects.png}
    \caption{Framing effects on alignment scores. The safety-first framing consistently yields higher alignment than neutral or freedom-first framings.}
    \label{fig:framing}
\end{figure}

% Example 3: Two figures side-by-side
\begin{figure}[H]
    \centering
    \begin{subfigure}[b]{0.48\textwidth}
        \centering
        \includegraphics[width=\textwidth]{figures/judge_agreement_overall.png}
        \caption{Overall agreement}
        \label{fig:judge_overall}
    \end{subfigure}
    \hfill
    \begin{subfigure}[b]{0.48\textwidth}
        \centering
        \includegraphics[width=\textwidth]{figures/judge_agreement_by_temp.png}
        \caption{By temperature}
        \label{fig:judge_temp}
    \end{subfigure}
    \caption{Judge agreement patterns. (a) Overall agreement rates per model. (b) Agreement rates by temperature and model.}
    \label{fig:judge}
\end{figure}

% Example 4: Three figures in a row
\begin{figure}[H]
    \centering
    \begin{subfigure}[b]{0.32\textwidth}
        \centering
        \includegraphics[width=\textwidth]{figures/temp_gemma.png}
        \caption{gemma:2b}
        \label{fig:temp_gemma}
    \end{subfigure}
    \hfill
    \begin{subfigure}[b]{0.32\textwidth}
        \centering
        \includegraphics[width=\textwidth]{figures/temp_llama.png}
        \caption{llama3:8b}
        \label{fig:temp_llama}
    \end{subfigure}
    \hfill
    \begin{subfigure}[b]{0.32\textwidth}
        \centering
        \includegraphics[width=\textwidth]{figures/temp_mistral.png}
        \caption{mistral:7b}
        \label{fig:temp_mistral}
    \end{subfigure}
    \caption{Temperature effects for three models.}
    \label{fig:temp_models}
\end{figure}

% Example 5: Figure with reference in text
As shown in Figure \ref{fig:temperature}, alignment scores vary modestly across temperatures. 
The range of variation (0.15--0.23 points) indicates relative stability, with orca-mini:7b 
showing the highest overall alignment (mean 4.33).

% ============================================
% FIGURE SIZING OPTIONS
% ============================================

% Full width
\includegraphics[width=\textwidth]{figures/filename.png}

% 80% width (recommended for most figures)
\includegraphics[width=0.8\textwidth]{figures/filename.png}

% Fixed width
\includegraphics[width=10cm]{figures/filename.png}

% Fixed height
\includegraphics[height=8cm]{figures/filename.png}

% Scale factor
\includegraphics[scale=0.8]{figures/filename.png}

% Both width and height (may distort)
\includegraphics[width=10cm, height=8cm]{figures/filename.png}

% ============================================
% CHECKLIST
% ============================================

% Before submitting to arXiv:
% [ ] All PNG files have no spaces or special characters in names
% [ ] File names in \includegraphics{} match exactly (case-sensitive)
% [ ] \usepackage{graphicx} is in the preamble
% [ ] PDFLaTeX compiler is selected (not regular LaTeX)
% [ ] figures/ directory is included in upload
% [ ] Tested compilation locally
% [ ] All figures have captions and labels
% [ ] References use \ref{fig:label} syntax

